% !TeX spellcheck = en_US
% !TeX encoding = UTF-8
\chapter*{Abstract}

In the third lab of the Security Insider Lab, we train machine learning (ML) and deep learning (DL) models to detect DDoS flows in the Friday-afternoon file of CIC-IDS2017 (225,745 flows). After removing infinite and missing values, duplicate rows and constant and identical columns, 223,082 flows with 59 features remain. We split them 70/15/15 (stratified) into train, validation and test sets and standardize the features with statistics from the training set only. We train five classical classifiers (Gaussian Naive Bayes, Logistic Regression, K-Nearest Neighbors, Random Forest, XGBoost) and three neural networks in PyTorch (MLP, LSTM, 1D-CNN), and compare all eight on the 33,463 test flows. All models reach F1 scores of 0.9909 or higher. XGBoost makes no error at all, Random Forest makes 3 errors and K-Nearest Neighbors 6. The best neural network is the MLP (F1 0.9995, 18 errors), and Gaussian Naive Bayes is the weakest (F1 0.9909) because of 327 false alarms. The LSTM first trained unstably, and gradient clipping with a tuned learning rate and patience reduced the effect on the result (F1 0.9943 to 0.9992), although the validation loss still oscillates. Because the scores are so close, we compare false-positive and false-negative counts and zoomed ROC curves instead of accuracy alone. The near-perfect results partly rest on a shortcut feature, the destination port. In Tasks 4 and 5 we survey eBPF/XDP, write an XDP per-source packet-rate limiter and build and test xdp-tools: with a limit of 100 packets/s per source, 697,064 flood packets sent by one attacker led to only 500 reaching the victim, while benign traffic was unaffected.
